% Part: many-valued-logic
% Chapter: syntax-and-semantics
% Section: introduction

\documentclass[../../../include/open-logic-section]{subfiles}

\begin{document}

\olfileid{mvl}{syn}{int}

\olsection{引言}

在经典逻辑中，我们处理由命题变元通过命题联结词
$\lnot$、$\land$、$\lor$、$\lif$ 和 $\liff$ 构建而成的公式。
当我们为经典逻辑定义语义时，我们使用两个真值 $\True$ 和 $\False$。
我们通过赋值~$\pAssign{v}$ 来解释命题变元，该赋值将真值 $\True$ 和 $\False$ 指派给命题变元。
任何赋值都为任意公式~$!A$ 确定一个真值 $\pValue{v}(!A)$，
并且一个公式在赋值~$\pAssign{v}$ 中被满足，记作 $\pSat{v}{!A}$，当且仅当 $\pValue{v}(!A) = \True$。

多值逻辑是经典二值逻辑的推广，它允许使用多于 $\True$ 和 $\False$ 的真值。
因此，在多值逻辑中，赋值~$\pAssign{v}$ 是一个函数，它从一系列可能的真值中为每个命题变元~$p$ 指派一个真值。
我们通常将允许的真值集合记作~$V$。
经典逻辑是 $V = \{\True, \False\}$ 的多值逻辑，
并且真值~$\pValue{v}(!A)$ 是使用我们所熟知的、联结词的特征真值表来计算的。

一旦增加了额外的真值，如何计算我们读作“与”、“或”、“非”和“如果——那么”的联结词的~$\pValue{v}(!A)$，就有了不止一种自然的选择。
因此，一个多值逻辑不仅由真值集合决定，还由我们为每个联结词选定使用的\emph{真值函数}决定。
然而，一旦为某个多值逻辑~$\Log L$ 选定了这些要素，真值 $\pValue{v}(!A)[\Log L]$ 便像经典逻辑中一样，由赋值唯一确定。
与经典逻辑一样，多值逻辑也是\emph{真值函数的}。

有了这些语义构件，我们就可以进一步定义重言式、语义后承和可满足性等语义概念的类似概念。
在经典逻辑中，如果对任何~$\pAssign{v}$，公式的真值 $\pValue{v}(!A) = \True$，则该公式是重言式。
在多值逻辑中，我们也需要对此作一点推广。
首先，真值集合~$V$ 并不要求包含~$\True$。
例如，一些多值逻辑使用数字作为真值集合，比如所有介于 $0$ 和~$1$ 之间的有理数。
在这种情况下，$1$~通常扮演~$\True$ 的角色。
在其他逻辑中，扮演这一角色的可能不止一个真值，而是多个真值。
因此，我们要求每个多值逻辑都有一个\emph{指定值}的集合~$V^+$。
据此，我们可以说一个公式在赋值~$\pAssign{v}$ 中被满足，记作 $\pSat{v}{!A}[\Log L]$，当且仅当 $\pValue{v}(!A)[\Log L] \in V^+$。
一个公式~$!A$ 是逻辑~$\Log L$ 的重言式，记作 $\Entails[\Log L] !A$，当且仅当对任何~$\pAssign{v}$，$\pValue{v}(!A) \in V^+$。
最后，如果每个满足~$\Gamma$ 中所有公式的赋值也都满足~$!A$，我们就说 $!A$ 是公式集合~$\Gamma$ 的语义后承，记作 $\Gamma \Entails[\Log L] !A$。

\end{document}